\par\Needspace{4\baselineskip}
\begin{center}
\hypertarget{native-056-article038}{}\pdfbookmark[2]{38.}{article-038}38.
\end{center}
It will be useful to illustrate the preceding article with an example. Let the true anomaly be $=100^0$ and the eccentricity $=0{,}967\,6456\,7$, with $\log q=9{,}765\,6500$. Here is the calculation of $E$, $B$, $A$, and $t$.
\begin{center}
\begin{tabular}{@{}ll@{}}
\makebox[18em][l]{$\log\operatorname{tang}\tfrac{1}{2}v$\hfill .\,.\,.\,.} & $0{,}076\,1865$\\[.3em]
\makebox[18em][l]{$\log\BareRoot{\dfrac{1-e}{1+e}}$\hfill .\,.\,.\,.\,.} & $9{,}107\,9927$\\[.4em]
\hline
\makebox[18em][l]{$\log\operatorname{tang}\tfrac{1}{2}E$\hfill .\,.\,.} & $9{,}184\,1792$,
\end{tabular}
\end{center}
\noindent This gives $\tfrac{1}{2}E=8^0\,41^{\prime}\,19^{\prime\prime}{,}32$ and $E=17^0\,22^{\prime}\,38^{\prime\prime}{,}64$. To this value of $E$ corresponds $\log B=0{,}000\,0040$. Further, in parts of the radius we find $E=0{,}303\,2928$ and $\sin E=0{,}298\,6643$. This gives $\tfrac{9}{20}E+\tfrac{1}{20}\sin E=0{,}151\,4150$, whose logarithm is $=9{,}180\,1689$, and consequently $\log A^{\frac{1}{2}}=9{,}180\,1649$. We obtain the following by formula 1) of the preceding article.
\begin{center}
\begin{tabular}{@{}lrl l@{}}
\multicolumn{3}{@{}l}{\makebox[22em][l]{$\log\dfrac{2Bq^{\frac{3}{2}}}{k\BareRoot{(1-e)}}$\hfill .\,.\,.\,.\,.\,.\,.\,.\,.}} & $2{,}458\,9614$\\[1.0em]
\multicolumn{3}{@{}l}{\makebox[22em][l]{$\log A^{\frac{1}{2}}$\hfill .\,.\,.\,.\,.\,.\,.\,.\,.\,.\,.\,.}} & $9{,}180\,1649$\\[.3em]
\hline
$\log$ & $43{,}56386$ & $=$ & $1{,}639\,1263$\\[.3em]
 & $19{,}98014$ & & \\
\cline{2-2}
 & $63{,}54400$ & $=t$ &
\end{tabular}
\end{center}
\begin{center}
\begin{tabular}{@{}ll@{}}
\makebox[22em][l]{$\log\dfrac{2B(1+9e)}{15k}\left(\dfrac{q}{1-e}\right)^{\frac{3}{2}}$\hfill .\,.\,.} & $3{,}760\,1038$\\[.4em]
\makebox[22em][l]{$\log A^{\frac{3}{2}}$\hfill .\,.\,.\,.\,.\,.\,.\,.\,.\,.\,.\,.} & $7{,}540\,4947$\\[.3em]
\hline
$\log 19{,}98014$\hfill $=$ & $1{,}300\,5985$
\end{tabular}
\end{center}
\noindent Working the same example by the usual method, we find $e\sin E$ in \SourceLeaf{057}arcseconds to be $=59610^{\prime\prime}{,}79=16^0\,33^{\prime}\,30^{\prime\prime}{,}79$, giving a mean anomaly $=49^{\prime}\,7^{\prime\prime}{,}85=2947^{\prime\prime}{,}85$. From this and $\log k\left(\dfrac{1-e}{q}\right)^{\frac{3}{2}}=1{,}666\,4302$ we obtain $t=63{,}54410$. The difference, here only a $10\,000^{\mathrm{th}}$ of a day, could easily have been three or four times as large if the errors had acted in the same direction.

It is also clear that a table of this kind for $\log B$ suffices by itself to solve the inverse problem with any desired precision. We determine $E$ by repeated trials so that the resulting value of $t$ agrees with the given one. This procedure would be rather laborious, so we shall now show how to arrange the auxiliary table much more conveniently, dispense with unguided trials altogether, and reduce the entire calculation to a neat and rapid algorithm that seems to leave nothing to be desired.
